% SPDX-License-Identifier: CC-BY-4.0
\documentclass[11pt,letterpaper]{article}

\usepackage[T1]{fontenc}
\usepackage{amsmath,amssymb}
\usepackage{tikz}
\usepackage{pgfplots}
\pgfplotsset{compat=1.18}
\usepackage{algorithm}
\usepackage{algpseudocode}
\usepackage{booktabs}
\usepackage{enumitem}
\usepackage[hidelinks]{hyperref}

\newcommand{\N}{\mathbb{N}}
\newcommand{\Z}{\mathbb{Z}}
\newcommand{\Q}{\mathbb{Q}}
\newcommand{\R}{\mathbb{R}}
\newcommand{\C}{\mathbb{C}}
\newcommand{\Map}[3]{#1 \colon #2 \to #3}
\newcommand{\Mod}[1]{#1\text{-}\mathrm{Mod}}
\newcommand{\GroupHom}{\mathrm{Hom}_{\mathbf{Grp}}}
\newcommand{\RingHom}{\mathrm{Hom}_{\mathbf{Ring}}}

\newtheorem{theorem}{Theorem}
\newtheorem{lemma}[theorem]{Lemma}
\newtheorem{corollary}[theorem]{Corollary}
\newtheorem{definition}[theorem]{Definition}
\newtheorem{remark}[theorem]{Remark}
\newenvironment{proof}{\noindent\textit{Proof.}}{\hfill\(\square\)}

\title{A finite invariant core}
\author{ascendnoosphere}
\date{\today}

\begin{document}
\maketitle

\begin{abstract}
Let \(A\), \(B\) and \(m\) be natural numbers with \(A<m\), and write
\(N=\lfloor B/(m-A)\rfloor\). If a map \(H\colon\Z\to\Z\) satisfies
\(m\,|H(n)|\le A\,|n|+B\) at every integer, then \(|n|\ge N+1\) forces
\(|H(n)|<|n|\), while \(|n|\le N\) forces \(|H(n)|\le N\). Every orbit
therefore enters the finite invariant interval of radius \(N\) and is
eventually periodic there. The argument does not construct a first-division
return, and it does not decide a classical convergence statement.
\end{abstract}

\section{The radius}

Throughout, \(A,B,m\in\N\) and \(A<m\). In particular \(m\ge1\) and the
positive gap \(m-A\) is a legal divisor. Absolute values of integers are
identified with their values in \(\N\).

\begin{definition}\label{Def:Radius}
The core radius is
\[
N=\left\lfloor\frac{B}{m-A}\right\rfloor .
\]
\end{definition}

\begin{lemma}\label{Lem:Room}
\(B<(m-A)(N+1)\).
\end{lemma}

\begin{proof}
Write \(d=m-A\), so \(d\ge1\) and \(B=dN+(B\bmod d)\) with
\(0\le B\bmod d<d\). A natural number strictly below \(d\) is at most
\(d-1\), and \(d-1<d\), hence
\[
B\le dN+(d-1)<dN+d=d(N+1),
\]
which is the claim.
\end{proof}

\begin{lemma}\label{Lem:Descent}
If \(n\ge N+1\), then \(A n+B<m n\).
\end{lemma}

\begin{proof}
Lemma~\ref{Lem:Room} and \(N+1\le n\) give
\(B<d(N+1)\le d n\). Therefore
\[
A n+B<A n+d n=A n+(m n-A n)=m n,
\]
where the rearrangement \(d n=m n-A n\) uses \(A\le m\).
\end{proof}

\begin{lemma}\label{Lem:Invariant}
If \(n\le N\), then \(A n+B<m(N+1)\).
\end{lemma}

\begin{proof}
The inequality \(n\le N\) gives \(A n\le A N\), so
\(A n+B\le A N+B\). Lemma~\ref{Lem:Room} then yields
\[
A N+B<A N+d(N+1).
\]
Expanding the second term,
\begin{align*}
A N+d(N+1)
&=A N+d N+d
=(A+d)N+d
=m N+d.
\end{align*}
Since \(d=m-A\le m\),
\[
m N+d\le m N+m=m(N+1).
\]
The inequality before this last comparison is strict, so
\(A n+B<m(N+1)\).
\end{proof}

\section{Maps obeying the bound}

\begin{definition}\label{Def:Bound}
A map \(H\colon\Z\to\Z\) obeys the core bound when
\[
m\,|H(k)|\le A\,|k|+B
\]
for every \(k\in\Z\).
\end{definition}

No identification of \(H\) with a division of an affine block is used below.
The only input is Definition~\ref{Def:Bound}.

\begin{theorem}\label{Thm:Map}
Suppose \(H\) obeys the core bound. If \(|n|\ge N+1\), then
\(|H(n)|<|n|\). If \(|n|\le N\), then \(|H(n)|\le N\).
\end{theorem}

\begin{proof}
In the first case Lemma~\ref{Lem:Descent} supplies
\(A\,|n|+B<m\,|n|\). Combined with the bound,
\(m\,|H(n)|<m\,|n|\). Cancelling the positive factor \(m\) gives
\(|H(n)|<|n|\). In the second case Lemma~\ref{Lem:Invariant} supplies
\(A\,|n|+B<m(N+1)\), so \(m\,|H(n)|<m(N+1)\) and therefore
\(|H(n)|<N+1\). Since both sides are natural numbers,
\(|H(n)|\le N\).
\end{proof}

\begin{theorem}\label{Thm:Entry}
If \(H\) obeys the core bound, then for every \(n\in\Z\) some iterate
\(H^{j}(n)\) satisfies \(|H^{j}(n)|\le N\).
\end{theorem}

\begin{proof}
Proceed by induction on \(|n|\). If \(|n|\le N\), the zero-th iterate is
already inside the radius. If \(|n|\ge N+1\), Theorem~\ref{Thm:Map} gives
\(|H(n)|<|n|\). The inductive hypothesis applied to \(H(n)\) produces an
entry time \(j\), and \(j+1\) is an entry time for \(n\).
\end{proof}

\begin{lemma}\label{Lem:Stays}
If \(H\) obeys the core bound and \(|n|\le N\), then
\(|H^{j}(n)|\le N\) for every \(j\).
\end{lemma}

\begin{proof}
The case \(j=0\) is the hypothesis. If the \(j\)-th iterate lies in the
radius, Theorem~\ref{Thm:Map} places its successor in the radius as well.
\end{proof}

\begin{theorem}\label{Thm:Periodic}
If \(H\) obeys the core bound, then every orbit is eventually periodic:
there exist \(t\ge0\) and \(p\ge1\) such that
\[
H^{t+k}(n)=H^{t+k+p}(n)
\]
for every \(k\ge0\).
\end{theorem}

\begin{proof}
Theorem~\ref{Thm:Entry} gives a time \(t\) at which the state
\(x=H^{t}(n)\) satisfies \(|x|\le N\). Lemma~\ref{Lem:Stays} keeps every
later state in the same interval, which contains \(2N+1\) integers. Among
the \(2N+2\) states \(x,H(x),\dots,H^{2N+1}(x)\), two therefore agree, say
\(H^{i}(x)=H^{j}(x)\) with \(i<j\). The gap \(p=j-i\) is positive. Shifting
the equality by an arbitrary \(k\) gives
\(H^{i+k}(x)=H^{j+k}(x)\), because both sides are the \(k\)-th iterate of
the same state. Reassembling the clocks,
\[
H^{t+i+k}(n)=H^{i+k}(x)=H^{j+k}(x)=H^{t+j+k}(n),
\]
and \(t+j+k=t+i+k+p\).
\end{proof}

\begin{remark}\label{Rem:Scope}
The floor quotient \(n\mapsto\lfloor(A\,|n|+B)/m\rfloor\), decorated with a
sign, obeys Definition~\ref{Def:Bound}, so Theorems~\ref{Thm:Map},
\ref{Thm:Entry} and \ref{Thm:Periodic} apply to it. That quotient is not
claimed to be a first-division return of a raw map. The hitting time of a
zero residue, the residue-dependent length of an affine block, the choice of
\(A\) and \(B\) as maxima over residues, traps, and every classical
convergence statement remain outside this argument.
\end{remark}

\begin{table}[htbp]
\centering
\begin{tabular}{ll}
\toprule
Statement & Formal name \\
\midrule
Definition~\ref{Def:Radius} & \texttt{CoreRadius} \\
Lemma~\ref{Lem:Room} & \texttt{CoreBoundRoom} \\
Lemma~\ref{Lem:Descent} & \texttt{CoreDescentNat} \\
Lemma~\ref{Lem:Invariant} & \texttt{CoreInvariantNat} \\
Definition~\ref{Def:Bound} & hypothesis of the map theorems \\
Theorem~\ref{Thm:Map} & \texttt{CoreMapDescent}, \texttt{CoreMapInvariant} \\
Theorem~\ref{Thm:Entry} & \texttt{CoreEntry} \\
Lemma~\ref{Lem:Stays} & \texttt{CoreStays} \\
Theorem~\ref{Thm:Periodic} & \texttt{CoreEventuallyPeriodic} \\
Remark~\ref{Rem:Scope} & non-formal scope summary \\
\bottomrule
\end{tabular}
\caption{Formal correspondence of the statements.}
\label{Table:SignedCoreLean}
\end{table}

\end{document}
